\subsection{Conditional Independence}


\subsubsection{Independence for Random Variables}

\begin{Motivation}
    If we throw two dice with outcome values $X$ and $Y$, resp., then knowing the value of $Y$ does not give us any information about the value of $X$, and vice versa. We say that $X$ and $Y$ are independent from each other. We will formalize this intuition for all random variables in the following.
\end{Motivation}

\begin{Def}[Independence of two random variables]
    \label{def:indep-rv}
    Let $(\Wcal,P(W))$ be a probability space and $X:\, \Wcal \to \Xcal$ and $Y:\, \Wcal \to \Ycal$ be two random variables.
    We say that $X$ and $Y$ are \emph{independent} if the following equation holds:
    \[
        P(X,Y) = P(X) \otimes P(Y),
    \]
    where $P(X,Y)$ is the joint and $P(X)$ and $P(Y)$ are the corresponding marginal distributions.
    In symbols we would write this as:
    \[ X \Indep_{P(W)} Y.\]
\end{Def}

\begin{Lem}
    \label{lem:indep-rv}
     Let $(\Wcal,P(W))$ be a probability space, $\Xcal$ and $\Ycal$ standard measurable spaces 
     and $X:\, \Wcal \to \Xcal$ and $Y:\, \Wcal \to \Ycal$ be two random variables.
     Then the following statements are equivalent:
     \begin{enumerate}
         \item $\displaystyle X \Indep_{P(W)} Y$.
         \item $P(X,Y) = P(X) \otimes P(Y)$.
         \item There exists a probability distribution $Q(X)$ such that:
             \[ P(X,Y) = Q(X) \otimes P(Y).\]
         \item $P(X|Y) = P(X)$ holds $P(Y)$-almost-surely, 
             where $P(X|Y)$ is a version of a conditional probability distribution from Cor.\ \ref{cor:reg-cond-prob-distr}. 
         \item For all $A \in \Bcal_\Xcal$ we have:
             \[ \E[\I_A(X)|Y] = \E[\I_A(X)] \qquad P(W)\text{-a.s.}\]
         \item For all $A \in \Bcal_\Xcal$ and $B \in \Bcal_\Ycal$ we have:
             \[ P(X \in A, Y \in B) = P(X \in A) \cdot P(Y \in B). \]
             %\[ \E[\I_A(X) \cdot \I_B(Y)] = \E[\I_A(X)] \cdot \E[\I_B(Y)]. \]
     \end{enumerate}
     \begin{proof}
         Exercise.
     \end{proof}
\end{Lem}

\begin{Exc}
    Reformulate the statements in Lemma \ref{lem:indep-rv} for the case we either have mass functions (discrete case) 
    or densities w.r.t.\ a product measure, e.g.\ the Lebesgue measure (absolute continuous case).
\end{Exc}

%\subsubsection{Mutual Independence of Random Variables}

We can generalize the notion of independence to arbitrary families of random variables:

\begin{Def}[Mutual independence of families of random variables]
    Let $(\Wcal,P(W))$ be a probability space and $I$ an (index) set. For $i \in I$ let $X_i:\, \Wcal \to \Xcal_i$ be 
    a random variable. We say that $(X_i)_{i \in I}$ is \emph{(mutually/jointly) independent}  if
    for all two disjoint subsets $J_1 \dcup J_2 \ins I$ we have the independence:
    \[ (X_{j_1})_{j_1 \in J_1} \Indep_{P(W)} (X_{j_2})_{j_2 \in J_2}.\]
\end{Def}

\begin{Exc}[Mutual independence for finite tuples of random variables]
    A finite tuple of random variables $(X_1,\dots,X_n)$ is mutually independent if and only if:
    \[ P(X_1,\dots,X_n) = P(X_1) \otimes \cdots \otimes P(X_n).\]
\end{Exc}

\begin{comment}
\begin{Rem}[Mutual independence for families of random variables]
%    \begin{enumerate}
%        \item 
    If $I$ is an arbitrary\footnote{Note that the theorem of Ionescu-Tulcea works for the \emph{independent} product of probability measures for \emph{arbitrary} index sets, not just for \emph{countable} index sets. So in this case, one does not need to invoke Kolmogorov's extension theorem, see \cite{Kle20} Thm.\ 14.39, which requires standard measurable spaces, but, which also works for dependent products for arbitrary index sets.} index set one can define the product measure: $\bigotimes_{ i \in I} P(X_i)$, by the theorem of Ionescu-Tulcea, see \cite{Tul49, Lam87} and theorem \ref{thm:ionescu-tulcea}, 
    which is determined by its finite marginals. 
    One can then show that $(X_i)_{i \in I}$ is mutually independent if and only if:
    \[ P\lp (X_i)_{i \in I} \rp = \bigotimes_{ i \in I} P(X_i).\]
%\item If $I$ is an arbitrary (not necessarily countable) index set and $\Xcal_i$ are standard measurable spaces for $i \in I$ then one can define the product measure: $\bigotimes_{ i \in I} P(X_i)$, by Kolmogorov's extension theorem, see \cite{Kle20} Thm.\ 14.39, which is already determined by its finite marginals: $P(X_{i_1}) \otimes \cdots \otimes P(X_{i_n})$, $i_1,\dots,i_n \in I$, $n \in \N$. With these preliminaries, we can show that $(X_i)_{i \in I}$ is mutually independent if and only if:
%    \[ P\lp (X_i)_{i \in I} \rp = \bigotimes_{ i \in I} P(X_i).\]
%    \end{enumerate}
%theorem of Ionescu-Tulcea \cite{Tul49, Lam87,Kle20} Thm.\ 14.35 and Cor.\ 14.36
\end{Rem}
\end{comment}

\begin{Exc}
    Let $(\Wcal,P(W))$ be a probability space and $I$ an arbitrary index set. 
    For $i \in I$ let $X_i:\, \Wcal \to \Xcal_i$ be a random variable. 
    The following statements are equivalent:
    \begin{enumerate}
    \item  $(X_i)_{i \in I}$ is (mutually/jointly) independent.
    \item For every finite disjoint subsets $J_1,J_2 \ins I$ we have the independence:
        \[ (X_{j_1})_{j_1 \in J_1} \Indep_{P(W)} (X_{j_2})_{j_2 \in J_2}.\]
    \item For every finite subset $J \ins I$ and $i \in I \sm J$ we have:
           \[ X_i \Indep_{P(W)} (X_j)_{j \in J}.\]
    \item For every finite subset $J \ins I$ we have:
           \[ P\lp (X_j)_{j \in J} \rp = \bigotimes_{j \in J} P(X_j).\]
    \item We have the equality:
            \[ P\lp (X_i)_{i \in I} \rp = \bigotimes_{ i \in I} P(X_i),\]
            where $\bigotimes_{ i \in I} P(X_i)$ is the product measure on $\Xcal = \prod_{ i \in I} \Xcal_i$, which is determined by the corresponding products on its finite marginals via the extension theorem of Ionescu-Tulcea, see \cite{Tul49, Lam87} and theorem \ref{thm:ionescu-tulcea}.
    \end{enumerate}
\end{Exc}


\subsubsection{Conditional Independence for Random Variables}


\begin{Motivation}
    \begin{enumerate}
        \item Consider two independent coin flips with outcome variables $X$ and $Y$, resp., with values in $\{0,1\}$, and
    $Z:=X + Y \in \{0,1,2 \}$. If the value of $Z$ is known, say $Z=1$, then revealing the value of $Y$, say $Y=0$, provides us with all the information to fully determine the value of $X$, here $X=0$. This is despite the fact that $X$ and $Y$ were assumed to be independent. This means that conditioning on a third variable $Z$ can destroy independence. In this case, we say that $X$ and $Y$ are dependent conditioned on $Z$. Summarized in symbols we have:
    \[ X \Indep_{P(W)} Y, \qquad \text{but:} \qquad X \nIndep_{P(W)} Y \given Z.\]
\item Now consider three (mutually) independent coin flips $X$, $W$, $U$ with values in $\{0,1\}$.
    Let $Z:=X +W $ and $Y:=Z+U = X + W + U$. If we knew the value of $Y$, say $Y=0$, then we would have information about the values of $X$ as well, here $X=0$. This shows that $X$ and $Y$ can not be independent random variables.
    If, in contrast, we would first reveal the value of $Z$, say $Z=1$, then the value of $X$ might be restricted by the value of $Z$, but also revealing $Y$ would not give us any \emph{additional} information about the value of $X$. The reason is that $Y = Z + U$ and $U$ is independent of $X$, $W$ and $Z=X+W$. So, even though $X$ and $Y$ are dependent, when conditioned on $Z$ they become independent, as there is no additional information gained about each others value, when revealing the other. Summarized in symbols we have:
        \[ X \nIndep_{P(W)} Y, \qquad \text{but:} \qquad X \Indep_{P(W)} Y \given Z.\]
    \end{enumerate}
    We now want to formalize conditional independence for random variables.
\end{Motivation}

\begin{Rem}[Conditional independence]
    In contrast to (unconditional) independence, see Definition \ref{def:indep-rv}, possible definitions of conditional independence come with many more subtleties, due to their interplay with conditional probability distributions or conditional expectations. Such definitions can in general be non-equivalent. However, if we restrict ourselves to standard measurable spaces the subtleties can be resolved and the definitions become equivalent. 
    This is the reason that in the following we will only state conditional independence for standard measurable spaces.
    We will make a clearer choice later for conditional independence of conditional random variables.
\end{Rem}

\begin{DefLem}[Conditional independence for random variables]
    \label{deflem:cond-indep-rv}
        Let $(\Wcal,P(W))$ be a probability space and $\Xcal$, $\Ycal$ and $\Zcal$ standard measurable spaces, and $X:\, \Wcal \to \Xcal$ and $Y:\, \Wcal \to \Ycal$ and $Z:\, \Wcal \to \Zcal$ be three random variables. 
        We then say that \emph{$X$ is independent of $Y$ conditioned on $Z$} if any of the following equivalent conditions holds:
    \begin{enumerate}
        \item $\displaystyle P(X,Y|Z) = P(X|Z) \otimes P(Y|Z)$ holds  $P(Z)$-a.s.,
            where $P(X,Y|Z)$, $P(X|Z)$, $P(Y|Z)$ are versions of conditional probability distributions from Cor.\ \ref{cor:reg-cond-prob-distr}.
        \item $\displaystyle P(X|Y,Z) = P(X|Z)$ holds  $P(Y,Z)$-a.s., where $P(X|Y,Z)$, $P(X|Z)$ from Cor.\ \ref{cor:reg-cond-prob-distr}.
        \item There exists a Markov kernel $Q(X|Z):\, \Zcal \dshto \Xcal$ such that:
            \[ P(X,Y,Z) = Q(X|Z) \otimes P(Y,Z).\]
        \item For every $A \in \Bcal_\Xcal$ and $B \in \Bcal_\Ycal$ we have:
            \[ \E\lB \I_A(X) \cdot \I_B(Y) |Z\rB = \E\lB \I_A(X) |Z\rB \cdot \E\lB \I_B(Y) |Z\rB\qquad P(W)\text{-a.s.} \]        \item For every $A \in \Bcal_\Xcal$ we have:
            \[ \E\lB \I_A(X)  |Y, Z\rB = \E\lB \I_A(X) |Z\rB \qquad P(W)\text{-a.s.} \]
    \end{enumerate}
    In those cases, in symbols we write:
    \[ X \Indep_{P(W)} Y \given Z.\]
    \begin{proof}
        Exercise.
    \end{proof}
\end{DefLem}

\begin{Exc}
    Restate all the statements in Definition/Lemma \ref{deflem:cond-indep-rv} for discrete random variables in terms of mass functions.
\end{Exc}


Conditional independence for random variables satisfies the following rules:

\begin{Thm}[Separoid axioms for conditional independence for random variables, see \cite{Daw01}]
    \label{thm:sym-sep-ax-cond-ind}
  Let $(\Wcal,P(W))$ be a probability space and $X$, $Y$, $Z$ and $U$ random variables taking values in standard measurable spaces $\Xcal, \Ycal, \Zcal$ and $\Ucal$, respectively.
    Then we have the following rules:
    \begin{enumerate}
        \item Redundancy: If $U= \varphi(X)$ a.s. is a measurable function of $X$, e.g.\ $U=X$, then:
            \[ U \Indep_{P(W)} Y \given X. \]
        \item Symmetry: 
            \[ X \Indep_{P(W)} Y \given Z \qquad \implies \qquad Y \Indep_{P(W)} X \given Z.\]
        \item Decomposition: 
            \[X \Indep_{P(W)} Y,U \given Z \qquad \implies \qquad X \Indep_{P(W)} U \given Z.\]
        \item Weak Union:
            \[X \Indep_{P(W)} Y,U \given Z \qquad \implies \qquad X \Indep_{P(W)} Y \given U, Z.\]
        \item Contraction:
            \[ \lp X \Indep_{P(W)} U \given Z  \rp \quad \land \quad \lp X \Indep_{P(W)} Y \given U,Z  \rp \qquad \implies \qquad X \Indep_{P(W)} Y,U \given Z.  \]
    \end{enumerate}
    \begin{proof}
        Exercise. 
    \end{proof}
\end{Thm}

\begin{Rem}
    The separoid axioms, see Theorem \ref{thm:sym-sep-ax-cond-ind}, also hold true for random variables that map into general (non-standard) measurable spaces if 
    one restricts oneself to the definition of conditional independence only involving conditional expectations (rather than conditional probabilities or Markov kernels) in Definition \ref{deflem:cond-indep-rv}.
\end{Rem}

\begin{Exc}
    Assume that the random variables $X$, $Y$, $Z$, $U$ have a joint density $p$ w.r.t.\ some product measure (or a joint mass function)
    such that $p(y,u|z) > 0$ for all values $y,u,z$. Show that we then also have the following \emph{intersection} rule:
            \[ \lp X \Indep_{P(W)} U \given Y,Z  \rp \quad \land \quad \lp X \Indep_{P(W)} Y \given U,Z  \rp \qquad \implies \qquad X \Indep_{P(W)} Y,U \given Z.  \]

\end{Exc}


\subsubsection{Conditional Independence for Conditional Random Variables}

\begin{Motivation}
    Assume that we are given a statistical model $P(X|\Theta)$ and a statistic $S=S(X)$, which is a measurable function of $X$.
    Often one wants to find such an $S$ such that the choice of parameter $\Theta=\theta$ has no ``influence'' on the probability distribution of $X$ when $S$ is provided. Such a statistic is usually called a \emph{sufficient} statistic of $X$ w.r.t.\ $P(X|\Theta)$. In symbols we want $S$ such that:
    \[ X \Indep \Theta \given S.\]
    However, the parameter variable $\Theta$ here is not a proper random variable as we have no distribution $P(\Theta)$ specified over it. Still such a conditional independence statement makes sense. We thus want to formalize a notion of conditional independence for conditional random variables. We follow the definition of \cite{For21}.
    Other approaches can be found in \cite{Daw79, Daw80, Daw01, CD17, RERS17, FM20}.
\end{Motivation}

\begin{Motivation}
    Consider a probabilistic program with input variables $T$, $S$ and output variables $X$, $Y$, $Z$. 
    Whenever the program is given $T$ and $S$ as input, it internally samples $U,E \sim U[0,1]$ uniformly and independently from a random number generator,  then calculates:
    \begin{align*}
        X &:= T + S + U, & Y &:= 5 \cdot S + E, & Z&:= X \cdot Y,
    \end{align*}
    and, finally, outputs $X$, $Y$ and $Z$.
    Even though, the input $T$ and $S$ is provided by the user and is not considered a random variable, we can reason about the fact that 
    ``Output $Y$ only depends on the input $S$ and not on the input $T$.'' 
    We want to formalize such conditional indpendence mathematically in order to be able to write this as:
    \[ Y \Indep T \given S.\]
\end{Motivation}

\begin{Def}[Conditional independence for \trv{}s]
\label{cond-indep-markov-kernels-def}
Let $(\Wcal \times \Tcal, K(W|T))$ be a transition probability space with Markov kernel:
\[K(W|T): \, \Tcal \dashrightarrow  \Wcal.\]
Consider conditional random variables: 
\[ X:\, \Wcal \times \Tcal \to \Xcal, \qquad Y:\,\Wcal \times \Tcal \to \Ycal, \qquad
Z:\,\Wcal \times \Tcal \to \Zcal.\]
We say that \emph{$X$ is independent of $Y$ conditioned on $Z$ w.r.t.\ $K(W|T)$}, in symbols:
\[X \Indep_{K(W|T)} Y \given Z, \]
if there \underline{\emph{exists} a Markov kernel}: 
\[ Q(X|Z):\; \Zcal \dashrightarrow \Xcal,\]
such that:
\[ K(X,Y,Z|T) = Q(X|Z) \otimes K(Y,Z|T),\]
where $K(Y,Z|T)$ is the marginal of $K(X,Y,Z|T)$.\footnote{
For the equation $K(X,Y,Z|T) = Q(X|Z) \otimes K(Y,Z|T)$ to hold it is sufficient to check that for all $t \in \Tcal$, $A \in \Bcal_\Xcal$, $B \in \Bcal_\Ycal$ and $C \in \Bcal_\Zcal$ we have:
\[ K(X \in A, Y \in B, Z \in C |T=t) = \int_C \int_B  Q(X \in A|Z=z)\, K(Y \in dy, Z \in dz|T=t). \]
}\\
As a special case, we define:
\[X \Indep_{K(W|T)} Y \qquad :\iff \qquad X \Indep_{K(W|T)} Y \given \ast. \]
\end{Def}



\begin{Not}[Essential uniqueness]
        The Markov kernel $Q(X|Z)$ appearing in the conditional independence $X \Indep_{K(W|T)} Y \given Z$ in definition \ref{cond-indep-markov-kernels-def} is then a version of a conditional Markov kernel $K(X|Y,Z,T)$ and is thus essentially unique in the sense of \ref{ess-unique}.
    We will use the following suggestive notation for it:
    \[ K(X|\cancel{T,Y},Z) := Q(X|Z),\]
    or similarly with crossed variables in different order.
    So we have in case of $X \Indep_{K(W|T)} Y \given Z$:
\[ K(X,Y,Z|T) = K(X|\cancel{T,Y},Z)\otimes K(Y,Z|T).\]
Note that $K(X|\cancel{T,Y},Z)$ is a version of the conditional Markov kernel $K(X|Y,Z,T)$ and does not depend on arguments $y$ and $t$.
\end{Not}

\begin{Rem}[Conditional independence includes conditional independence from $T$]
    \label{add-T}
 We have the equivalence:
 \[X \Indep_{K(W|T) } Y \given Z \quad\iff \quad X \Indep_{K(W|T)} T,Y \given Z,\]
where $T :\, \Wcal \times \Tcal \to \Tcal$, $(w,t) \mapsto t$, is the canonical projection map.
\begin{proof}
\begin{align*}
&X \Indep_{K(W|T)} Y \given Z \\
\iff &\, \exists Q(X|Z):\, K(X,Y,Z|T) &&= Q(X|Z) \otimes K(Y,Z|T)\\
\iff &\, \exists Q(X|Z):\, K(X,T,Y,Z|T) &&= Q(X|Z)\otimes \underbrace{ K(Y,Z|T)  \otimes  \delta(T|T)  }_{K(T,Y,Z|T)}\\
\iff& X \Indep_{K(W|T)} T,Y \given Z.
\end{align*}
The middle implication ``$\Longrightarrow$'' follows by taking the product with $\delta(T|T)$, 
and the reverse implication ``$\Longleftarrow$'' by marginalizing out $T$, i.e.\ via $\delta(T \in \Tcal|T)=1$.
\end{proof}
\end{Rem}


\begin{Rem}[How to find $Q(X|Z)$ and check for conditional independence?]
 In case we have the conditional independence:
 \[X \Indep_{K(W|T)} Y \given Z, \]
  we then get by definition:
 \[ K(X,Y,Z|T) = Q(X|Z)\otimes K(Y,Z|T),\]
 for some Markov kernel $Q(X|Z)$. This implies for all $t \in \Tcal$ the equation:
\[ K(X,Z|T=t) = Q(X|Z)\otimes K(Z|T=t).\]
This means that  $Q(X|Z)$ is a version of the conditional probability distribution $K(X|Z,T=t)$ 
for all $t \in \Tcal$ at once, and, in addition, it is also functionally not dependent on $t$. 
So for fixed $t_0 \in \Tcal$ the conditional $K(X|Z,T=t_0)$ can be changed on a $K(Z|T=t_0)$-null-set such that it agrees with $Q(X|Z)$.
So it is reasonable to test out versions of $K(X|Z,T=t_0)$ for $Q(X|Z)$. 
To summarize, we have the following equivalence between:
\begin{enumerate}
    \item $\displaystyle X \Indep_{K(W|T)} Y \given Z$,
    \item There exist $t_0 \in \Tcal$ and a (regular) version of the conditional probability distribution 
        $K(X|Z,T=t_0)$ such that for all $t \in \Tcal$: 
            \[ K(X,Y,Z|T=t) = K(X|Z,T=t_0)\otimes K(Y,Z|T=t). \]
       Note that in the last expression the middle term has the fixed $t_0$ and the outer two terms have varying $t \in \Tcal$.
\end{enumerate}
This equivalence allows us to narrow our search space for $Q(X|Z)$ to such conditional probability distributions.
\end{Rem}


\begin{Eg}[Conditional independence for discrete conditional random variables]
    Let the situation be like in definition \ref{cond-indep-markov-kernels-def} and assume all spaces to be countable and discrete.
    Let $k$ be the mass function for $K(X,Y,Z|T)$.
    Then we have: 
    \[X \Indep_{K(W|T)} Y \given Z, \]
    if and only if there is a probability mass function $q$ such that for all values $x,y,z,t$:
    \[ k(x,y,z|t) = q(x|z) \cdot k(y,z|t). \]
    Note that in this case $q(x|z)$ is a version of $k(x|z,t)$ for all $t \in \Tcal$ at once, but that is also independent of $t$.
    We can use this knowledge to find such a proposal $q(x|z)$ as follows.

    If there exists a $t_0 \in \Tcal$ such that $k(z|t_0) >0$ for all  $z \in \Zcal$ then the 
    conditional $k(x|z,t_0)$ is uniquely given and equal to $\frac{k(x,z|t_0)}{k(z|t_0)}$.
    $k(x|z,t_0)$ would then necessarily agree with $q(x|z)$ in case of the conditional independence.
    So, if there exists a $t_0 \in \Tcal$ such that $k(z|t_0) >0$ for all $z \in \Zcal$ then 
    we get the following equivalence:
    \[X \Indep_{K(W|T)} Y \given Z \qquad \iff \qquad \forall x,y,z,t: \qquad k(x,y,z|t) = k(x|z,t_0) \cdot k(y,z|t). \]
    Again, note that in the last expression the middle term has the fixed $t_0$ and the outer two terms have varying $t \in \Tcal$.
\end{Eg}

This example can be generalized.

\begin{Thm}[Conditional independence for conditional random variables with density]
    Let $\mu_\Xcal$, $\mu_\Ycal$ and $\mu_\Zcal$ reference measures on $\Xcal$, $\Ycal$ and $\Zcal$, resp., and $\mu:=\mu_\Xcal \otimes \mu_\Ycal \otimes \mu_\Zcal$ the product measure on
    $\Xcal\times\Ycal\times\Zcal$. Assume that $K(X,Y,Z|T)$ has a Doob-Radon-Nikodym derivative $k$ w.r.t.\ $\mu$ and let $t_0 \in \Tcal$ be a fixed value.
    Then we have the implication:
    \[\forall t \in \Tcal\, \forall_\mu\, x,y,z. \qquad k(x,y,z|t) = k(x|z,t_0) \cdot k(y,z|t) \qquad \implies \qquad X \Indep_{K(W|T)} Y \given Z, \]
    where $\forall_\mu$ means ``for $\mu$-almost-all''. If for $\mu_\Zcal$-almost-all $z \in \Zcal$ we have: $k(z|t_0)>0$, then also the reverse implication holds, with $k(x|z,t_0) := \frac{k(x,z|t_0)}{k(z|t_0)}$:
    \[\forall t \in \Tcal\, \forall_\mu\, x,y,z. \qquad k(x,y,z|t) = k(x|z,t_0) \cdot k(y,z|t) \qquad \Longleftarrow \qquad X \Indep_{K(W|T)} Y \given Z. \]
    \begin{proof}
        ``$\implies$'': This is clear, as the factorization of the densities provides the needed factorization of the corresponding Markov kernels.\\
    ``$\Longleftarrow$'': By assumption we have a factorization:
    \[ K(X,Y,Z|T) = Q(X|Z) \otimes K(Y,Z|T),  \]
    which implies for all $A \in \Bcal_\Xcal$ and $C \in \Bcal_\Zcal$:
    \begin{align*} 
        &\int_C \int_A k(x,z|t_0) \, \mu_\Xcal(dx)\, \mu_\Zcal(dz) \\
        &=  K(X \in A,Z \in C|T=t_0) \\ 
        &= \int_C Q(X \in A|Z=z)\,K(Z \in dz|T=t_0)  \\ 
        &= \int_C Q(X \in A|Z=z)\,k(z|t_0)\, \mu_\Zcal(dz). 
    \end{align*}
    This implies that we have:
    \begin{align*} 
        \forall_{\mu_\Zcal}\, z \in \Zcal.\qquad   \int_A k(x,z|t_0) \, \mu_\Xcal(dx)  =  Q(X \in A|Z=z)\,k(z|t_0), 
    \end{align*}
    which implies, since $k(z|t_0) > 0$, that $k(x|z,t_0) = \frac{k(x,z|t_0)}{k(z|t_0)}$ is a density of $Q(X|Z)$ up to a $\mu_\Zcal$-null set $N$.
    Since $K(Z|T) \ll \mu_\Zcal$ this $N$ is also a $K(Z|T)$-null set, and thus $\Ycal \times N$ a $K(Y,Z|T)$-null set. 
    So for all $t \in \Tcal$, $A \in \Bcal_\Xcal$, $B \in \Bcal_\Ycal$, $C \in \Bcal_\Zcal$ we get:
        \begin{align*} 
          &  \int_C \int_B \int_A k(x,y,z|t) \, \mu_\Xcal(dx)\, \mu_\Ycal(dy)\, \mu_\Zcal(dz) \\
        &=  K(X \in A, Y \in B, Z \in C|T=t) \\ 
        &= \int_C \int_B Q(X \in A|Z=z)\,K(Y \in dy, Z \in dz|T=t)  \\ 
        &= \int_C \int_B \int_A k(x|z,t_0) \,\mu_\Xcal(dx)\, K(Y \in dy, Z \in dz|T=t)  \\ 
        &= \int_C \int_B \int_A k(x|z,t_0) \,\mu_\Xcal(dx)\, k(y,z|t) \, \mu_\Ycal(dy)\, \mu_\Zcal(dz)  \\ 
        &= \int_C \int_B \int_A k(x|z,t_0) \cdot k(y,z|t) \,\mu_\Xcal(dx) \, \mu_\Ycal(dy)\, \mu_\Zcal(dz).
    \end{align*}
    So the corresponding Markov kernels on the lhs and rhs are the same.
    This implies that the set:
    \[ M:= \lC (x,y,z,t) \st k(x,y,z|t) \neq k(x|z,t_0) \cdot k(y,z|t) \rC \]
    is a $\mu$-null set in $\Bcal_\Xcal\otimes\Bcal_\Ycal\otimes\Bcal_\Zcal\otimes\Bcal_\Tcal$.
    So for all $t \in \Tcal$ and $\mu$-almost-all $x,y,z$ the following equation holds:
    \[k(x,y,z|t) = k(x|z,t_0) \cdot k(y,z|t), \]
    which implies the claim.
    \end{proof}
\end{Thm}


\begin{Rem}[Conditional independence for random variables]
    \label{ci-random-variables}
    By Definition/Lemma \ref{deflem:cond-indep-rv} we recover the notion of conditional independence for random variables 
    $X,Y,Z$ with standard measurable spaces as codomains by taking $\Tcal:=\{\ast\}$, $P(W):=K(W|\ast)$:
\[ X \Indep_{P(W)} Y \given Z \qquad \iff \qquad \exists Q(X|Z):\, P(X,Y,Z) = Q(X|Z) \otimes P(Y,Z).\]
Such a $Q(X|Z)$ is then a conditional probability distribution of $P(X,Z)$ conditioned on $Z$.
In suggestive notations:
\[Q(X|Z) =:P(X|\cancel{Y},Z) = P(X|Z).\]
\end{Rem}



\begin{Lem}[Conditional independence for deterministic mappings]
    Let $F:\, \Tcal \to \Fcal$ and $H:\,\Tcal \to \Hcal$ be measurable mappings, with $\Fcal$ standard.
    We now consider them as (deterministic) \trv{}s on the transition probability space $(\Wcal \times \Tcal,K(W|T))$ via:
    \[\begin{array}{ccccc} 
        F: & \Wcal  \times \Tcal &\to &\Fcal,\\
           &(w,t) &\mapsto& F(t), \\
        H: & \Wcal  \times \Tcal &\to &\Hcal,\\
           &(w,t) &\mapsto& H(t), \\
       % G: & \Wcal  \times \Tcal &\to &\Gcal,\\
       %    &(w,t) &\mapsto& G(t), \\
    \end{array}\]
    which do not depend on the 'probabilistic part' $\Wcal$ of $K(W|T)$. Let $G:\, \Wcal \times \Tcal \to \Gcal$ be another \trv{}.\\
    We write $F \precsim H$ if there exists a measurable map $\varphi:\,\Hcal \to \Fcal$ such that:
    \[ F = \varphi \circ H. \]
    Then we have the equivalence:\footnote{The full equivalence needs Kuratowski's extension theorem for standard measurable spaces
    (see \cite{Kec95} 12.2): Any measurable map from a (not necessarily measurable) subset of a measurable space to a standard measurable space extends to a measurable map on the whole space. Alternatively, one could define $F \precsim H$ via existence of measurable $\varphi:\, H(\Tcal) \to \Fcal$ such that $F=\varphi \circ H$, but this moves problems elsewhere.}
    \[ F \precsim H \qquad \iff\qquad F \Indep_{K(W|T)} G \given H.\]
    So $F$ is a deterministic measurable map of $H$ iff $F$ is independent of $G$ given $H$. Note that the first part of the statement is independent of $G$.% and we as well take $G=T$.
\begin{proof}
 ``$\implies$'': This direction is rather easy. See the later separoid axioms.\\
 ``$\Longleftarrow$'': Since $F$ and $H$ are deterministic and only dependent on $T$ we get that:
    \[ K(F,G,H|T) = \delta(F|T) \otimes \delta(H|T) \otimes K(G|T).  \]
    By the conditional independence we now have a Markov kernel $Q(F|H)$ such that we have the factorization:
    \[ K(F,G,H|T) = Q(F|H) \otimes K(G,H|T) = Q(F|H) \otimes \delta(H|T) \otimes K(G|T).  \]
    Marginalizing out $G$, $H$ and taking $T=t$ we get from these equations:
    \[ \delta_{F(t)} = \delta(F|T=t) = Q(F|H(t)),\]
    which is a Dirac measure centered at $F(t)$.
    We can now define the mapping:
    \[\varphi:\, H(\Tcal) \to \Fcal, \quad H(t) \mapsto F(t), \]
    which is well-defined, because $h:=H(t_1)=H(t_2)$ implies that $Q(F|H=h)$ is a Dirac measure centered at $F(t_1)$ and $F(t_2)$, so $F(t_1)=F(t_2)$.
    $\varphi$ is measurable.
    Indeed, its composition with $\delta:\,\Fcal \to \Pcal(\Fcal), \, z \mapsto \delta_z$ equals $Q(F|H)$, which is measurable.
    Since $\Bcal_\Fcal = \delta^*\Bcal_{\Pcal(\Fcal)}$, see lemma \ref{lem-maps-markov-kernel} 2., also $\varphi$ is measurable.
    Since $\Fcal$ is a standard measurable space, $\varphi$ extends to a measurable mapping $\varphi:\, \Hcal \to \Fcal$ by Kuratowski's extension theorem for standard measurable spaces (see \cite{Kec95} 12.2). 
    Finally, note that we have $F(t) = \varphi(H(t))$ for all $(w,t) \in \Wcal \times \Tcal$, which shows the claim.
\end{proof}
\end{Lem}


\begin{Eg}[Conditional independence for deterministic mappings]
    If for example, $\Tcal = \Tcal_1 \times \Tcal_2$ and $T_i:\,\Wcal \times \Tcal_1 \times \Tcal_2 \to \Tcal_i$ the canonical projection onto $\Tcal_i$,
    then $F$ is a function in two variables $(t_1,t_2)$.
    We then  have:
    \[ F \Indep_{K(W|T)} T_1 \given T_2,\]
    if and only if $F$---as a function---is only dependent on the argument $t_2$ (and not on $t_1$).
    %So $F$ is then independent of all the (other) arguments $t=(t_1,t_2)$ given/besides $t_2$.
\end{Eg}

Another example of what conditional independence of \trv{}s can encode is the following.

\begin{Rem}[Existence of conditional Markov kernels expressed as conditional independence]
    \label{conditional-markov-kernel-as-ci}
    Let $X$, $Y$ be conditional random variables on transition probability space $(\Wcal \times \Tcal, K(W|T))$.
    Then we can express the existence of a conditional Markov kernel $K(X|Y,T)$ as the conditional independence:
    \[ X\Indep_{K(W|T)} \ast|Y,T,  \]
    where $\ast$ is the constant \trv{}. Alternatively and equivalently, we could also write:
    \[ X\Indep_{K(W|T)} T|Y,T.  \]
    Note that for standard measurable spaces $\Xcal$ and $\Ycal$ the above statement always holds.
    In suggestive symbols:
    \[ K(X|\cancel{T},Y,T) = K(X|Y,T). \]
\end{Rem}

\begin{Eg}[Certain statistics expressed as conditional independence] Let $P(W|\Theta)$ be a statistical model, considered as a Markov kernel $\Fcal \dshto \Wcal$. Let $X$ and $Y$ be two \trv{}s w.r.t.\ $P(W|\Theta)$.
    A \emph{statistic} of $X$ is a measurable map $S:\,\Xcal \to \Scal$, which we consider as the  \trv{} $S\precsim X$ given via:
    \[ S:\,\Wcal \times \Fcal \to \Scal,\quad (w,\theta) \mapsto S(X(w,\theta)). \]
\begin{enumerate}
\item \emph{Ancillarity}. $S$ is an \emph{ancillary statistic} of $X$ w.r.t.\ $\Theta$ if and only if:
\[ S \Indep_{P(W|\Theta)} \Theta.\]
This means that every parameter $\Theta=\theta$ induces the same distribution for $S$: 
\[P(S|\Theta=\theta) = P(S|\cancel{\Theta}).\]
\item \emph{Sufficiency}. $S$ is a \emph{sufficient statistic} of $X$ w.r.t.\ $\Theta$ if and only if:
\[ X \Indep_{P(W|\Theta)} \Theta \given S.\]
This means that there is a Markov kernel $P(X|S,\cancel{\Theta})$ such that:
\[ P(X,S|\Theta) = P(X|S,\cancel{\Theta}) \otimes P(S|\Theta).\]
So $X$ only ``interacts'' with the parameters $\Theta$ through $S$.
%\item[] A measurable map $\Phi: \Tcal \to \Fcal$ is called \emph{sufficient parameterization} for $X$ w.r.t.\ $P(W|T)$ if and only if:
%\[ X \Indep_{P(W|T),\exists} T \given \Phi.\]
%For example, $\Phi$ defined by $\Phi(t):=P(X|T=t) \in \Pcal(\Xcal)=:\Fcal$ is always a sufficient parameterization for $X$ (by the same arguments as in example \ref{propensity}).
\item \emph{Adequacy}. $S$ is an \emph{adequate statistic} of $X$ for $Y$ w.r.t.\ $\Theta$ if and only if:
\[ X \Indep_{P(W|\Theta)} \Theta,Y \given S.\]
This means we have a factorization:
\[ P(X,Y,S|\Theta) = P(X|\cancel{\Theta,Y},S) \otimes P(Y,S|\Theta),\]
for some Markov kernel $P(X|\cancel{\Theta,Y},S)$.
This means that all information of $X$ about the parameters and labels $Y$ is fully captured already by $S$.
\end{enumerate}
\end{Eg}

%\begin{comment}
\begin{Thm}
    \label{thm-ci-equivalences}
Let $(\Wcal \times \Tcal, K(W|T))$ be a transition probability space with Markov kernel:
\[K(W|T): \, \Tcal \dashrightarrow  \Wcal.\]
Consider \trv{}s $X$, $Y$, $Z$ with common domain $\Wcal\times \Tcal$ and codomains $\Xcal$, 
$\Ycal$ and $\Zcal$, resp., and $T:\,\Wcal\times \Tcal \to \Tcal$ the canonical projection map.
We will write $P(X|Z)=K(X|\cancel{T},Z)$ 
for a fixed version of the Markov kernel appearing in the conditional independence $X \Indep_{K(W|T)} T \given Z$ (only in case it holds).\\
With these notations, the following are equivalent:
\begin{enumerate}
  \item $\displaystyle X \Indep_{K(W|T)} Y \given Z,$ \label{thm-ci-equivalences:base}
    \item $\displaystyle X \Indep_{K(W|T)} T, Y \given Z,$
    \item  $\displaystyle X \Indep_{K(W|T)} T \given Z$ and  $ \displaystyle K(X,Y,Z|T) = P(X|Z) \otimes K(Y,Z|T)$.
    \item  $\displaystyle  X \Indep_{K(W|T)} T \given Z$ and for every $t \in \Tcal$ we have:
      $\displaystyle X_t \Indep_{K(W|T=t)} Y_t \given Z_t.$ \label{thm-ci-equivalences:every-t}
    % \item  $\displaystyle  X \Indep_{K(W|T)} T \given Z$ and for every $t \in \Tcal$ and $A \in \Bcal_\Xcal$ we have:
    %     \[\displaystyle \E_t[\I_A(X_t)|Y_t,Z_t] = \E_t[\I_A(X_t)|Z_t] \qquad K(W|T=t)\text{-a.s.,}\]
    %     where the conditional expectation $\E_t$ is w.r.t.\ $K(W|T=t)$ (for each $t \in \Tcal$ individually).
\end{enumerate}
Furthermore, any of those points implies:
\[ K(X|\cancel{T,Y},Z) = K(X|\cancel{T},Z) \qquad K(Y,Z|T)\text{-a.s.}.\]
and the following:
\begin{enumerate}[resume]
    \item For every probability distribution $Q(T) \in \Pcal(\Tcal)$ we have the conditional independence\footnote{Note that this again implies the second part of point \ref{thm-ci-equivalences:every-t}: $X_t \Indep Y_t \given Z_t$ for every $t \in \Tcal$. So the first part of point \ref{thm-ci-equivalences:every-t}: $X \Indep T \given Z$, can then be seen as the additional obstruction to obtain the ``full'' conditional independence: $X \Indep Y \given Z$.}:
      \[ X \Indep_{K(W|T)\otimes Q(T) } T,Y \given Z.\] \label{thm-ci-equivalences:forall-Q}
\end{enumerate}
%\Joris{Perhaps explicitly state the obvious: Furthermore, 
%  \begin{enumerate}[resume]
%    \item $\displaystyle  X \Indep_{K(W|T)} T \given Z$ and \ref{thm-ci-equivalences:forall-Q}
%      implies \ref{thm-ci-equivalences:every-t}, and hence \ref{thm-ci-equivalences:base}.
%    This allows us to connect this with the notions in our UAI paper, and might give a general recipe to generalize probabilistic independence tests to the extended setting for conditional random variables.
%Note that $\displaystyle  X \Indep_{K(W|T)} T \given Z$ is equivalent to the existence of $Q$ s.t.\ $K(X,Z|T) = Q(X|Z)\otimes K(Z|T)$.
%    \item $\displaystyle  X \Indep_{K(W|T)} T \given Z$, and for every probability distribution $Q(T) \in \Pcal(\Tcal)$:
%        \[ X \Indep_{K(W|T)\otimes Q(T) } Y \given Z,T.\]
%        together implies that $\displaystyle X \Indep_{K(W|T)} Y,T \given Z.$
%        This is the slightly weaker version that Patrick wrote on my BB, which according to him matches even better with the UAI paper notions.
%    \item Yet another option would be: $\displaystyle  X \Indep_{K(W|T)} T \given Z$, and for every probability distribution $Q(T) \in \Pcal(\Tcal)$:
%        \[ X \Indep_{K(W|T)\otimes Q(T) } Y \given Z.\]
%        implies that $\displaystyle X \Indep_{K(W|T)} Y \given Z.$
%        This is obviously equivalent to 6, and the last part is most similar to definition C.1 in the UAI paper.
%  \end{enumerate}
%}
\begin{proof}
    3. $\implies$ 1. is clear by definition.\\
    1. $\iff$ 2.: by \ref{add-T}.\\
    2. $\implies$ 4.,5.:
    By assumption we have the factorization:
    \[K(X,Y,Z,T|T) = K(X|Z) \otimes K(Y,Z,T|T),\]
    for some Markov kernel $K(X|Z)$. Via marginalization and multiplication this implies the two equations:
    \begin{align*}
        K(X,Z,T|T)            &= K(X|Z) \otimes K(Z,T|T),\\
        \underbrace{K(X,Y,Z|T)\otimes Q(T)}_{=:Q(X,Y,Z,T)} &= K(X|Z) \otimes \underbrace{K(Y,Z|T) \otimes Q(T)}_{=Q(Y,Z,T)},
    \end{align*}
    for every $Q(T) \in \Pcal(\Tcal)$. The last equation shows 5.\\ 
    If we take $Q(T)=\delta_t$ we get:
    \[K(X_t,Y_t,Z_t|T=t) = K(X|Z_t) \otimes K(Y_t,Z_t|T=t). \]
    Together with the first of the above equations this shows 4.\\
%\begin{comment}
    4. $\implies$ 3.: By $X \Indep_{K(W|T)} T \given Z$ we have:
    \[K(X,Z|T) = P(X|Z) \otimes K(Z|T).\]
    By the assumption $ X_t \Indep_{K(W|T=t)} Y_t \given Z_t$, on the other hand, we have---for each $t \in \Tcal$ individually---a factorization:
    \[ K(X,Y,Z|T=t) = Q_t(X|Z) \otimes K(Y,Z|T=t),  \]
    with a Markov kernel $Q_t$, which might depend on $t \in \Tcal$, where we suppress the indices $t$ on all the variables for readability everywhere.
Marginalizing out $Y$ and comparing to the above we then get the two equalities:
\[ P(X|Z) \otimes K(Z|T=t) = K(X,Z|T=t)  = Q_t(X|Z) \otimes K(Z|T=t).\]
By the essential uniqueness of such a factorization we see that $P(X \in A|Z)$ only differs from $Q_t(X \in A|Z)$ on a $K(Z|T=t)$-null set. Considered as functions of $(y,z)$ (by ignoring $y$) they are equal up to a $K(Y,Z|T=t)$-null set. This means that we can replace $Q_t(X \in A|Z)$ with $P(X\in A|Z)$ for every $A \in \Bcal_\Xcal$ and $t \in \Tcal$. So we get the equation:
        \[ K(X,Y,Z|T=t) = P(X|Z) \otimes K(Y,Z|T=t),  \]
        for all $t \in \Tcal$ and thus:
  \[ K(X,Y,Z|T) = P(X|Z) \otimes K(Y,Z|T).  \]
  This shows 3.\\ 
%\end{comment}
\begin{comment}
%proof with even weaker assumptions, only using conditional expectations:
    \Joris{This proof relies on conditional expectations. Is that necessary? I guess that one could also find a direct proof that does not make use of the notion of conditional expectation. This would have my preference as it makes the lecture notes slightly more self-contained and accessible, and I find it refreshing that conditional independence is not defined via conditional expectations.} \PF{There was already an alternative proof in the tex file commented out. See now above. Below the other proof.}\\
    4. $\implies$ 3.:  By $X \Indep_{K(W|T)} T \given Z$ we have the factorization:
    \[K(X,Z|T) = P(X|Z) \otimes K(Z|T).\]
    This means that for every $t \in \Tcal$ and every measurable $A \ins \Xcal$, $C \ins \Zcal$ we have:
    \[ \E_t\lB \I_A(X_t) \cdot \I_C(Z_t)\rB = \E_t\lB P(X \in A|Z_t) \cdot \I_C(Z_t)\rB,  \]
    where the expectation $\E_t$ is w.r.t.\ $K(W|T=t)$. This shows that $P(X \in A|Z_t)$ is a version of $\E_t[\I_A(X_t)|Z_t]$ for every $t \in \Tcal$, 
    by the defining properties of conditional expectation.\\
    By the assumption $ X_t \Indep_{K(W|T=t)} Y_t \given Z_t$ we then have for every fixed $t\in \Tcal$ and measurable $A \ins \Xcal$:
    \[ \E_t[\I_A(X_t)|Y_t,Z_t] = \E_t[\I_A(X_t)|Z_t] = P(X \in A|Z_t)\qquad K(W|T=t)\text{-a.s.}\]
    By the defining properties of conditional expectation for $\E_t[\I_A(X_t)|Y_t,Z_t]$  we then get that for every measurable  $A \ins \Xcal$, $B \ins \Ycal$, $C \ins \Zcal$:
    \[ \E_t\lB \I_A(X_t) \cdot \I_B(Y_t) \cdot \I_C(Z_t)\rB = \E_t\lB P(X \in A|Z_t) \cdot \I_B(Y_t) \cdot \I_C(Z_t)\rB. \]
    Since this holds for every $t \in \Tcal$ we get:
    \[ K(X,Y,Z|T) = P(X|Z) \otimes K(Y,Z|T).\]
\end{comment}
\end{proof}
\end{Thm}
%\end{comment}


\begin{Rem}[Discrete $\Tcal$]
    In the setting of theorem \ref{thm-ci-equivalences},
    let  $\Xcal$, $\Zcal$ be standard measurable spaces and $\Tcal$ be a countable discrete measurable space with any fixed probability distribution $Q(T)$ that has a strictly positive mass function.
    Then we get the equivalence:
    \[X \Indep_{K(W|T)\otimes Q(T)} T,Y \given Z  \qquad \iff \qquad  X \Indep_{K(W|T)} Y \given Z.\]
    \begin{proof}
        The rhs implies the lhs side by theorem \ref{thm-ci-equivalences}. 
        So, now assume the lhs and put:
        \[ Q(X,Y,Z,T) := K(X,Y,Z|T) \otimes Q(T).  \]
        Its marginal is then denoted by $Q(X,Z)$.
        Since $\Xcal$, $\Zcal$ are standard measurable spaces we get a (regular) conditional probability distribution $Q(X|Z)$, such that $Q(X,Z) = Q(X|Z)\otimes Q(Z)$. By the assumed conditional independence we thus have:
        \begin{align*}
            K(X,Y,Z|T) \otimes Q(T) &= Q(X,Y,Z,T) \\
                                    &= Q(X|Z) \otimes Q(Y,Z,T) \\
                                    &= Q(X|Z) \otimes K(Y,Z|T) \otimes Q(T).
        \end{align*}
        Since $Q(T)$ is strictly positive and conditional Markov kernels are essentially unique we get the sure equality:
        \[K(X,Y,Z|T) = Q(X|Z) \otimes K(Y,Z|T),\]
        which implies the claim.
    \end{proof}
\end{Rem}

\begin{Cor}
%Let $(\Wcal \times \Tcal, K(W|T))$ be a transition probability space with Markov kernel:
%\[K(W|T): \, \Tcal \dashrightarrow  \Wcal.\]
%Consider \trv{}s $X$, $Y$, $Z$ with common domain $\Wcal\times \Tcal$ and codomains $\Xcal$, 
%$\Ycal$ and $\Zcal$, resp., where 
If $\Xcal$, $\Zcal$ are standard measurable spaces
then we have the equivalence:
\[ X \Indep_{K(W|T)} Y \given Z,T \qquad \iff \qquad \forall t \in \Tcal:\quad X_t \Indep_{K(W|T=t)} Y_t \given Z_t.\]
\begin{proof}
    This directly follows from theorem \ref{thm-ci-equivalences} 4.\ with $(Z,T)$ in the role of $Z$ 
    and remark \ref{conditional-markov-kernel-as-ci} to get the first part of 4.
    In suggestive symbols:
    \[ K(X|\cancel{T,Y},Z,T) = K(X|Z,T) \qquad K(Z|T)\text{-a.s.}\]
\end{proof}
\end{Cor}

\subsubsection{Example: Linear Gaussian Markov Kernels}

\begin{comment}
\begin{Lem}[Gaussian conditioning formula]
    If:
    \begin{align*}
        \matb{X\\Z} & \sim \Ncal\lp \matb{x\\z} \st \matb{\mu_X\\\mu_Z}, \matb{\Sigma_{X,X}&\Sigma_{X,Z}\\\Sigma_{Z,X}&\Sigma_{Z,Z}} \rp. 
    \end{align*}
    then:
    \begin{align*}
        (X|Z=z)  & \sim \Ncal\lp x \st \mu_{X|Z=z}, \Sigma_{X|Z} \rp, \\
    \mu_{X|Z=z} &:= \mu_X + \Sigma_{X,Z}\Sigma_{Z,Z}^{-1}\,(z-\mu_Z)\\
                &= \Sigma_{X,Z}\Sigma_{Z,Z}^{-1}\,z + \lp \mu_X - \Sigma_{X,Z}\Sigma_{Z,Z}^{-1}\,\mu_Z \rp, \\
    \Sigma_{X|Z} &:= \Sigma_{X,X} - \Sigma_{X,Z} \Sigma_{Z,Z}^{-1} \Sigma_{Z,X}. 
    \end{align*}
    Note that the latter covariance matrix does not depend on $z$.
\end{Lem}
\end{comment}


\begin{Thm}[Conditional independence for linear Gaussian conditional random variables]
Let $\Tcal = \RN$, $\Xcal = \RN^{d_X}$, $\Ycal = \RN^{d_Y}$ and $\Zcal = \RN^{d_Z}$. 
Consider a linear Gaussian Markov kernel $P(X,Y,Z|T)$, which is given by a density of the form:
\begin{align*}
    p(x,y,z|t) &= \Ncal\lp \matb{x\\y\\z} \st \matb{\Gamma_X\\\Gamma_Y\\\Gamma_Z} \cdot t + \matb{\gamma_X\\\gamma_Y\\\gamma_Z} , \matb{\Sigma_{X,X}&\Sigma_{X,Y}&\Sigma_{X,Z}\\\Sigma_{Y,X}&\Sigma_{Y,Y}&\Sigma_{Y,Z}\\\Sigma_{Z,X}&\Sigma_{Z,Y}&\Sigma_{Z,Z}} \rp. 
\end{align*}
%If $\det \Sigma_{Z,Z} \neq 0$,
Then we have the following equivalence:
\[X \Indep_{P(X,Y,Z|T)} Y \given Z \qquad \iff \qquad \Sigma_{X,Y} = \Sigma_{X,Z} \Sigma_{Z,Z}^{-1}\Sigma_{Z,Y} 
\quad \land \quad \Gamma_X = \Sigma_{X,Z}\Sigma_{Z,Z}^{-1}\,\Gamma_Z. \]
If this is the case then the Markov kernel $Q(X|Z)$ coming from the conditional independence:
\[P(X,Y,Z|T) = Q(X|Z) \otimes P(Y,Z|T),\]
is also a linear Gaussian Markov kernel with density:
\begin{align*}
    q(x|z) &= \Ncal\lp x \st \mu_{X|Z}(z), \Sigma_{X,X|Z} \rp,\\ 
    \mu_{X|Z}(z)   &:= \gamma_X + \Sigma_{X,Z}\Sigma_{Z,Z}^{-1}\, (z-\gamma_Z),\\
    \Sigma_{X,X|Z} & := \Sigma_{X,X} - \Sigma_{X,Z}\Sigma_{Z,Z}^{-1} \Sigma_{Z,X},
\end{align*}
which coincides with the usual marginal conditional for $t=0$, i.e.:
\[Q(X|Z=z) = P(X|Z=z,T=0).\]
So, we also get the equivalence:
\[X \Indep_{P(X,Y,Z|T)} Y \given Z \qquad \iff \qquad P(X,Y,Z|T) = P(X|Z,T=0) \otimes P(Y,Z|T). \]
\end{Thm}

\begin{proof}
First note that in general the conditional $P(X|Y,Z,T)$ is also a linear Gaussian Markov kernel and of the form:
\begin{align*}
p(x|y,z,t) &= \Ncal\lp x \st \mu_{X|Y,Z,T}(y,z,t), \Sigma_{X|Y,Z,T} \rp,
\end{align*}
with the following abbreviation for the covariance matrix:
\begin{align*}
    \Sigma_{X|Y,Z,T} & := \Sigma_{X,X} - \matb{\Sigma_{X,Y}&\Sigma_{X,Z}} \matb{\Sigma_{Y,Y}&\Sigma_{Y,Z}\\\Sigma_{Z,Y}&\Sigma_{Z,Z} }^{-1} \matb{\Sigma_{Y,X}\\\Sigma_{Z,X}},
\end{align*}
and the following abbreviation for the mean:
\begin{align*}
         &\mu_{X|Y,Z,T}(y,z,t) \\
         &:= \lp \Gamma_X \cdot t + \gamma_X \rp + \matb{\Sigma_{X,Y}&\Sigma_{X,Z}} \matb{\Sigma_{Y,Y}&\Sigma_{Y,Z}\\\Sigma_{Z,Y}&\Sigma_{Z,Z} }^{-1}\,\lp\matb{y\\z}-\matb{\Gamma_Y\cdot t + \gamma_Y\\\Gamma_Z \cdot t + \gamma_Z}\rp\\
                &= \matb{\Sigma_{X,Y}&\Sigma_{X,Z}} \matb{\Sigma_{Y,Y}&\Sigma_{Y,Z}\\\Sigma_{Z,Y}&\Sigma_{Z,Z} }^{-1}\,\matb{y\\z} \\
               &\qquad +  \lp \Gamma_X \cdot t + \gamma_X \rp - \matb{\Sigma_{X,Y}&\Sigma_{X,Z}} \matb{\Sigma_{Y,Y}&\Sigma_{Y,Z}\\\Sigma_{Z,Y}&\Sigma_{Z,Z} }^{-1}\,\matb{\Gamma_Y\cdot t + \gamma_Y\\\Gamma_Z \cdot t + \gamma_Z}.
    \end{align*}

We now want to investigate under which conditions we get the conditional independence:
\[ X \Indep_{P(X,Y,Z|T)} Y \given Z.\]
Note that in this case the conditional independence is equivalent to the statement that the conditional density $p(x|y,z,t)$
is not dependent on the arguments $y$ and $t$.

Let us first investigate the first term involving $y$:
\begin{align*}
    \matb{\Sigma_{X,Y}&\Sigma_{X,Z}} \matb{\Sigma_{Y,Y}&\Sigma_{Y,Z}\\\Sigma_{Z,Y}&\Sigma_{Z,Z}}^{-1}\,\matb{y\\z}.
\end{align*}
Note that we can use the following formula for the ($2 \times 2)$-block inverse:
\begin{align*}
  & \matb{\Sigma_{Y,Y}&\Sigma_{Y,Z}\\\Sigma_{Z,Y}&\Sigma_{Z,Z}}^{-1}  = \\
  & \matb{ \lp\Sigma_{Y,Y}-\Sigma_{Y,Z}\Sigma_{Z,Z}^{-1}\Sigma_{Z,Y}\rp^{-1} & 
    -\lp\Sigma_{Y,Y}-\Sigma_{Y,Z}\Sigma_{Z,Z}^{-1}\Sigma_{Z,Y}\rp^{-1} \Sigma_{Y,Z}\Sigma_{Z,Z}^{-1} \\
    -\Sigma_{Z,Z}^{-1}\Sigma_{Z,Y}\lp\Sigma_{Y,Y}-\Sigma_{Y,Z}\Sigma_{Z,Z}^{-1}\Sigma_{Z,Y}\rp^{-1} &
\Sigma_{Z,Z}^{-1} +\Sigma_{Z,Z}^{-1}\Sigma_{Z,Y}\lp\Sigma_{Y,Y}-\Sigma_{Y,Z}\Sigma_{Z,Z}^{-1}\Sigma_{Z,Y}\rp^{-1}\Sigma_{Y,Z}\Sigma_{Z,Z}^{-1} }
\end{align*}
This leads us to require that:
\begin{align*}
 \matb{\Sigma_{X,Y}&\Sigma_{X,Z}}   \matb{ \lp\Sigma_{Y,Y}-\Sigma_{Y,Z}\Sigma_{Z,Z}^{-1}\Sigma_{Z,Y}\rp^{-1} \\ 
    -\Sigma_{Z,Z}^{-1}\Sigma_{Z,Y}\lp\Sigma_{Y,Y}-\Sigma_{Y,Z}\Sigma_{Z,Z}^{-1}\Sigma_{Z,Y}\rp^{-1}} =0,
\end{align*}
which is equivalent to:
\begin{align*}
 \lp \Sigma_{X,Y} -\Sigma_{X,Z} \Sigma_{Z,Z}^{-1}\Sigma_{Z,Y} \rp \lp\Sigma_{Y,Y}-\Sigma_{Y,Z}\Sigma_{Z,Z}^{-1}\Sigma_{Z,Y}\rp^{-1}  =0,
\end{align*}
which can be further simplified, by multiplying with the inverse of the inverse, to:
\begin{align*}
     \Sigma_{X,Y} -\Sigma_{X,Z} \Sigma_{Z,Z}^{-1}\Sigma_{Z,Y}   = 0.
\end{align*}

We also need that the mean of the conditional is not dependent on $t$, which leads to the following condition coming from the second term:
\begin{align*}
    0 & = \Gamma_X -  \matb{\Sigma_{X,Y}&\Sigma_{X,Z}} \matb{\Sigma_{Y,Y}&\Sigma_{Y,Z}\\\Sigma_{Z,Y}&\Sigma_{Z,Z}}^{-1}\,\matb{\Gamma_Y\\\Gamma_Z} \\
    & =  \Gamma_X - \lp \Sigma_{X,Y} -\Sigma_{X,Z} \Sigma_{Z,Z}^{-1}\Sigma_{Z,Y} \rp \lp\Sigma_{Y,Y}-\Sigma_{Y,Z}\Sigma_{Z,Z}^{-1}\Sigma_{Z,Y}\rp^{-1} \, \Gamma_Y \\
    & \quad -  \Sigma_{X,Z}  \Sigma_{Z,Z}^{-1} \, \Gamma_Z 
 + \lp \Sigma_{X,Y} - \Sigma_{X,Z} \Sigma_{Z,Z}^{-1}\Sigma_{Z,Y}\rp\,\lp\Sigma_{Y,Y}-\Sigma_{Y,Z}\Sigma_{Z,Z}^{-1}\Sigma_{Z,Y}\rp^{-1} \Sigma_{Y,Z}\Sigma_{Z,Z}^{-1} \, \Gamma_Z \\
 &= \Gamma_X - \Sigma_{X,Z}\Sigma_{Z,Z}^{-1}\,\Gamma_Z,
\end{align*}
where we made repeated use of the condition: $\Sigma_{X,Y} - \Sigma_{X,Z} \Sigma_{Z,Z}^{-1}\Sigma_{Z,Y} =0$.

%So a necessary and sufficient condition for $X \Indep_{P(X,Y,Z|T)} Y \given Z$ is the following two equalities: 
%\begin{align*}
%    \Sigma_{X,Y} &= \Sigma_{X,Z} \Sigma_{Z,Z}^{-1}\Sigma_{Z,Y}, \\
%    \Gamma_X &= \Sigma_{X,Z}\Sigma_{Z,Z}^{-1}\,\Gamma_Z.
%\end{align*}

This leads us to the following equivalence for linear Gaussian Markov kernels:
\[X \Indep_{P(X,Y,Z|T)} Y \given Z \qquad \iff \qquad \Sigma_{X,Y} = \Sigma_{X,Z} \Sigma_{Z,Z}^{-1}\Sigma_{Z,Y} 
\quad \land \quad \Gamma_X = \Sigma_{X,Z}\Sigma_{Z,Z}^{-1}\,\Gamma_Z. \]
If this is the case then the Markov kernel $Q(X|Z)$ coming from the conditional independence:
\[P(X,Y,Z|T) = Q(X|Z) \otimes P(Y,Z|T),\]
is also a linear Gaussian Markov kernel
with density:
\begin{align*}
    q(x|z) &= \Ncal\lp x \st \mu_{X|Z}(z), \Sigma_{X,X|Z} \rp,\\ 
    \mu_{X|Z}(z)   &:= \gamma_X + \Sigma_{X,Z}\Sigma_{Z,Z}^{-1}\, (z-\gamma_Z),\\
    \Sigma_{X,X|Z} & := \Sigma_{X,X} - \Sigma_{X,Z}\Sigma_{Z,Z}^{-1} \Sigma_{Z,X},
\end{align*}
which is the usual marginal conditional for $t=0$.
\end{proof}



\subsection{Separoid Axioms for Conditional Independence}

The following asymmetric separiod axioms for conditional independence are a generalization of the symmetric separoid axioms due to A.P.~Dawid \cite{Daw01} and the similar graphoid axioms due to J.~Pearl and A.~Paz \cite{PP85}.
\begin{DefThm}[(Asymmetric) separoid axioms for conditional independence]
    \label{ci-separoid-axioms}
    Let $(\Wcal \times \Tcal, K(W|T))$ be a transition probability space with Markov kernel:
\[K(W|T): \, \Tcal \dashrightarrow  \Wcal.\]
Consider \trv{}s $X$, $Y$, $Z$, $U$ with common domain $\Wcal\times \Tcal$ and standard measurable spaces $\Xcal$, 
$\Ycal$, $\Zcal$, $\Ucal$, resp., as codomains.
Let $T:\, \Wcal \times \Tcal \to \Tcal$ be the canonical projection and $\ast$ the constant \trv{}.\\
We write $U \precsim X$ if there exists a measurable function %$G:\,X(\Wcal\times \Tcal) \to \Ucal$ 
$G:\, \Xcal \to \Ucal$ 
such that $U = G \circ X$.\\
%We write $X \lor Y$ to mean the \trv{} $(X,Y)$
Then the ternary relation $\Indep = \Indep_{K(W|T)}$ satisfies the following rules:
\begin{enumerate}[label=\alph*)]
	\item Extended Left Redundancy: 
    \item[] $U \precsim X \implies U \Indep  Y \given X$.
    \item $T$-Restricted Right Redundancy:\footnote{\label{fn:sep-ax-std} $T$-Restricted Right Redundancy, Left Weak Union and Symmetry need the existence of conditional Markov kernels. That is the reason we assumed standard measurable spaces.}
    \item[] $X\Indep \ast|Z,T$ always holds.
    \item $T$-Inverted Right Decomposition:
    \item[] $X \Indep  Y \given Z \implies X \Indep  T, Y \given Z$.
	\item Left Decomposition:  
	\item[]$X,U \Indep  Y \given Z \implies U \Indep  Y \given Z$.
  \item Right Decomposition:  	
    \item[] $X \Indep  Y,U \given Z \implies X \Indep  U \given Z$.
    \item Left Weak Union:\footref{fn:sep-ax-std}
    \item[] $X , U \Indep  Y  \given Z \implies X \Indep  Y \given U , Z$.
  \item Right Weak Union:\item[] $X \Indep  Y , U \given Z \implies X \Indep  Y \given U , Z$.
\item Left Contraction:\item[] 
 $ (X \Indep  Y \given U , Z) \land (U \Indep  Y \given Z) \implies X , U \Indep  Y \given Z$.
\item Right Contraction:
\item[] $ (X \Indep  Y \given U , Z) \land (X \Indep  U \given Z)  \implies X \Indep  Y , U \given Z$.
\item Right Cross Contraction:
\item[] $ (X \Indep  Y \given U , Z) \land (U \Indep  X \given Z) \implies X  \Indep  Y , U \given Z$.
\item Flipped Left Cross Contraction: 
\item[] $ (X \Indep  Y \given U , Z) \land (Y \Indep  U \given Z) \implies Y \Indep  X , U \given Z$.
\end{enumerate}
In particular, we have the equivalences:
$$ (X \Indep  Y , U \given Z) \quad\iff\quad (X \Indep  Y \given U , Z) \quad\land\quad (X \Indep  U \given Z)  ,$$
%and:
$$ (X , U \Indep  Y \given Z) \quad\iff\quad (X \Indep  Y \given U , Z) \quad\land\quad (U \Indep  Y \given Z).$$
We also get:
\begin{enumerate}[resume,label=\alph*)]
    \item $T$-Restricted Symmetry:\footref{fn:sep-ax-std} %If $\Tcal =\Asterisk$ then: 
	\item[] $X \Indep  Y \given Z,T \implies Y \Indep X \given Z,T$. 
\end{enumerate}
In the special case of $\Tcal = \Asterisk = \{\ast\}$, the one-point space, (i.e.\ in the case of probability distributions and random variables mapping to standard measurable spaces) we
thus have (unrestricted) Symmetry.
\end{DefThm}



